\section{Danh sách Ca Kiểm thử Chức năng Chi tiết Toàn diện}
Dưới đây là bảng đặc tả chi tiết của toàn bộ các ca kiểm thử chức năng được thiết kế cho hệ sinh thái Mioto.vn, chuẩn hóa theo chuẩn quốc tế IEEE 829:

\begin{longtable}{p{2.2cm}p{4cm}p{5.5cm}p{4cm}}
\caption{Danh mục chi tiết các ca kiểm thử chức năng hệ thống Mioto.vn} \label{tab:full_test_cases} \\
\toprule
\textbf{Mã TC} & \textbf{Tên ca kiểm thử} & \textbf{Các bước \& Dữ liệu (Steps)} & \textbf{Kỳ vọng (Expected Result)} \\
\midrule
\endfirsthead

\multicolumn{4}{c}{{\bfseries \tablename\ \thetable{} -- Tiếp tục từ trang trước}} \\
\toprule
\textbf{Mã TC} & \textbf{Tên ca kiểm thử} & \textbf{Các bước \& Dữ liệu (Steps)} & \textbf{Kỳ vọng (Expected Result)} \\
\midrule
\endhead

\midrule \multicolumn{4}{r}{{Còn tiếp ở trang sau...}} \\
\endfoot

\bottomrule
\endlastfoot

\textbf{TC-AUTH-01} & Đăng ký tài khoản với SĐT hợp lệ & 1. Truy cập trang Đăng ký. 2. Nhập SĐT hợp lệ. 3. Bấm nút "Tiếp tục". & Hệ thống thông báo gửi mã OTP thành công, chuyển sang màn hình nhập OTP, hiển thị đồng hồ đếm ngược 60 giây. \\ \midrule
\textbf{TC-AUTH-02} & Đăng ký với SĐT chứa ký tự chữ cái & 1. Nhập SĐT chứa chữ cái. 2. Bấm "Tiếp tục". & Nút "Tiếp tục" bị vô hiệu hóa hoặc báo lỗi: "Số điện thoại không đúng định dạng". \\ \midrule
\textbf{TC-AUTH-03} & Đăng ký với SĐT thiếu ký tự (9 chữ số) & 1. Nhập SĐT 9 số. 2. Bấm "Tiếp tục". & Báo lỗi: "Số điện thoại phải bao gồm đúng 10 chữ số". Không gửi OTP. \\ \midrule
\textbf{TC-AUTH-04} & Đăng ký với SĐT thừa ký tự (11 chữ số) & 1. Nhập SĐT 11 số. 2. Bấm "Tiếp tục". & Hệ thống tự động cắt ở ký tự thứ 10 hoặc báo lỗi không hợp lệ. \\ \midrule
\textbf{TC-AUTH-05} & Đăng ký với SĐT đã tồn tại trên hệ thống & 1. Nhập SĐT đã có tài khoản. 2. Bấm "Tiếp tục". & Thông báo lỗi: "Số điện thoại này đã được đăng ký. Vui lòng đăng nhập hoặc chọn Quên mật khẩu". \\ \midrule
\textbf{TC-AUTH-06} & Nhập đúng mã OTP trong thời hạn hiệu lực & 1. Nhập chính xác 6 chữ số OTP. 2. Bấm "Xác nhận". & Hệ thống xác thực thành công, chuyển sang bước "Thiết lập mật khẩu". \\ \midrule
\textbf{TC-AUTH-07} & Nhập sai mã OTP & 1. Nhập mã OTP không khớp. 2. Bấm "Xác nhận". & Báo lỗi: "Mã xác thực không chính xác. Vui lòng thử lại (Còn 4 lần thử)". \\ \midrule
\textbf{TC-AUTH-08} & Nhập mã OTP khi đã hết hạn (>60 giây) & 1. Chờ đồng hồ về 00:00. 2. Nhập mã OTP. 3. Bấm "Xác nhận". & Báo lỗi: "Mã OTP đã hết hiệu lực. Vui lòng bấm 'Gửi lại mã'". \\ \midrule
\textbf{TC-AUTH-09} & Nhập sai mã OTP liên tiếp 5 lần & 1. Nhập mã sai 5 lần liên tiếp. & Hệ thống khóa gửi/xác thực OTP trên SĐT này trong vòng 15 phút để chống Brute-force. \\ \midrule
\textbf{TC-AUTH-10} & Thiết lập mật khẩu thỏa mãn chính sách bảo mật & 1. Nhập mật khẩu đủ 8 ký tự, có hoa, thường, số, ký tự đặc biệt. 2. Nhập lại mật khẩu khớp. 3. Bấm "Hoàn tất". & Hệ thống thông báo đăng ký thành công, tự động đăng nhập và chuyển về Trang chủ. \\ \midrule
\textbf{TC-AUTH-11} & Thiết lập mật khẩu quá ngắn (< 8 ký tự) & 1. Nhập mật khẩu 5 ký tự. 2. Bấm "Hoàn tất". & Báo lỗi: "Mật khẩu phải chứa ít nhất 8 ký tự". Nút bấm bị vô hiệu hóa. \\ \midrule
\textbf{TC-AUTH-12} & Đăng nhập với SĐT và Mật khẩu chính xác & 1. Truy cập trang Đăng nhập. 2. Nhập SĐT và mật khẩu đúng. 3. Bấm "Đăng nhập". & Đăng nhập thành công, thanh tiêu đề hiển thị tên người dùng và ảnh đại diện, lưu JWT an toàn. \\ \midrule
\textbf{TC-AUTH-13} & Đăng nhập sai mật khẩu liên tiếp 5 lần & 1. Nhập sai mật khẩu liên tiếp 5 lần. & Hệ thống tạm khóa tài khoản trong 30 phút hoặc yêu cầu giải mã CAPTCHA để chống tấn công từ điển. \\ \midrule
\textbf{TC-AUTH-14} & Đăng nhập nhanh qua Google OAuth 2.0 & 1. Bấm nút "Đăng nhập với Google". 2. Chọn tài khoản Google trong popup. 3. Cho phép quyền truy cập cơ bản. & Đăng nhập thành công, hệ thống tự động liên kết email và chuyển về trang trước đó. \\ \midrule
\textbf{TC-AUTH-15} & Tải lên ảnh GPLX hợp lệ và gửi duyệt KYC & 1. Vào mục "Hồ sơ của tôi" -> "Giấy phép lái xe". 2. Tải 2 ảnh GPLX chuẩn. 3. Nhập số và ngày cấp. 4. Bấm "Gửi duyệt". & Hệ thống báo: "Hồ sơ GPLX của bạn đã được gửi thành công và đang chờ xét duyệt trong vòng 24h". Trạng thái đổi thành "Chờ duyệt". \\ \midrule
\textbf{TC-AUTH-16} & Tải lên ảnh bằng lái xe máy (Hạng A1/A2) & 1. Tải ảnh bằng lái xe máy A1. 2. Chọn hạng A1. 3. Bấm "Gửi duyệt". & Báo lỗi: "Chỉ chấp nhận Giấy phép lái xe ô tô hạng B1, B2, C, D, E, FC. Bằng lái xe máy không đủ điều kiện thuê ô tô". \\ \midrule
\textbf{TC-AUTH-17} & Nhập số GPLX đã hết hạn sử dụng & 1. Nhập thông tin bằng lái đã hết hạn. 2. Bấm "Gửi duyệt". & Báo lỗi: "Giấy phép lái xe đã hết hạn sử dụng. Vui lòng cập nhật GPLX còn hiệu lực để tiếp tục". \\ \midrule
\textbf{TC-SRCH-01} & Tìm kiếm xe theo Tỉnh/Thành phố hợp lệ & 1. Nhập "Hồ Chí Minh" vào ô tìm kiếm. 2. Chọn gợi ý "Thành phố Hồ Chí Minh, Việt Nam" từ danh sách autocomplete. 3. Bấm "Tìm xe". & Chuyển đến trang danh sách xe `/city/hochiminh` hoặc `/find/filter`, hiển thị toàn bộ xe đang hoạt động tại khu vực TP.HCM. \\ \midrule
\textbf{TC-SRCH-02} & Tìm kiếm xe với địa điểm không có trong danh mục & 1. Nhập địa danh không tồn tại. 2. Bấm "Tìm xe". & Danh sách gợi ý autocomplete rỗng, hiển thị thông báo: "Không tìm thấy địa điểm phù hợp". \\ \midrule
\textbf{TC-SRCH-03} & Chọn thời gian nhận xe trong quá khứ & 1. Mở lịch chọn ngày nhận. 2. Cố gắng chọn ngày trong quá khứ. & Các ngày trong quá khứ bị làm mờ (disabled) và không thể bấm chọn. \\ \midrule
\textbf{TC-SRCH-04} & Chọn thời gian nhận xe cách hiện tại dưới 2 giờ & 1. Chọn ngày hôm nay. 2. Chọn khung giờ 10:30. 3. Bấm "Tìm xe". & Báo lỗi hoặc tự động điều chỉnh: "Thời gian bắt đầu nhận xe phải sau thời điểm hiện tại ít nhất 2 giờ (tối thiểu 12:00)". \\ \midrule
\textbf{TC-SRCH-05} & Chọn thời gian thuê tròn 24 giờ (Hợp lệ) & 1. Chọn đúng khoảng cách 24h. 2. Bấm "Tìm xe". & Tìm kiếm thành công, hệ thống tính số ngày thuê là: 1 ngày. \\ \midrule
\textbf{TC-SRCH-06} & Chọn thời gian trả xe trước thời gian nhận xe & 1. Chọn ngày trả trước ngày nhận. 2. Bấm xác nhận. & Hệ thống ngăn chặn việc chọn, tự động đẩy ngày trả về sau ngày nhận ít nhất 1 ngày. \\ \midrule
\textbf{TC-SRCH-07} & Lọc xe theo loại hộp số: Tự động (Automatic) & 1. Tích chọn checkbox "Số tự động". 2. Quan sát danh sách xe. & 100\% các xe hiển thị trong danh sách đều có thông số hộp số là "Số tự động". Không có xe số sàn nào xuất hiện. \\ \midrule
\textbf{TC-SRCH-08} & Lọc kết hợp đa tiêu chí (Hãng xe + Hộp số + Nhiên liệu) & 1. Tích chọn hãng "VinFast". 2. Tích chọn "Số tự động". 3. Tích chọn nhiên liệu "Điện". & Danh sách chỉ trả về các dòng xe điện VinFast (VF e34, VF8, VF9, VF5). Số lượng xe khớp hiển thị chính xác trên tiêu đề. \\ \midrule
\textbf{TC-SRCH-09} & Lọc xe theo khoảng giá (Slider min - max) & 1. Kéo thanh slider giá min lên 800k, max xuống 1tr2. 2. Kiểm tra đơn giá các xe hiển thị. & Tất cả xe trả về đều có đơn giá nằm nghiêm ngặt trong khoảng từ 800.000đ đến 1.200.000đ/ngày. \\ \midrule
\textbf{TC-SRCH-10} & Lọc xe có tính năng "Đặt xe nhanh" (Instant Booking) & 1. Kích hoạt bộ lọc "Đặt xe nhanh". & Toàn bộ xe hiển thị đều có biểu tượng tia sét màu xanh (Instant Booking - Không cần chủ xe phê duyệt thủ công). \\ \midrule
\textbf{TC-SRCH-11} & Lọc xe có tiện ích "Miễn thế chấp tiền mặt / xe máy" & 1. Bật bộ lọc "Miễn thế chấp". & Danh sách chỉ hiển thị các xe có chính sách chấp nhận CCCD gắn chip thay thế cho việc giữ xe máy hoặc cọc 15 triệu tiền mặt. \\ \midrule
\textbf{TC-SRCH-12} & Đổi chế độ xem Danh sách sang Bản đồ số (Map View) & 1. Bấm vào nút chuyển đổi "Xem bản đồ". & Giao diện hiển thị bản đồ HERE Maps với các pin đánh dấu vị trí và giá thuê của từng xe tại tọa độ tương ứng. Khi di chuột/bấm vào pin sẽ hiển thị popup tóm tắt xe. \\ \midrule
\textbf{TC-PRIC-01} & Tính giá ngày thường, nhận tại bãi, không voucher (Quy tắc R1) & 1. Chọn xe giá 1.000.000đ. 2. Chọn thuê từ T3 đến T5 (2 ngày). 3. Xem bảng kê chi phí. & - Tiền thuê gốc: 2.000.000đ. - Phí bảo hiểm chuyến đi (10\%): 200.000đ. - Phí giao xe: 0đ. - Tổng cộng: 2.200.000đ. - Tiền cọc 30\%: 660.000đ. \\ \midrule
\textbf{TC-PRIC-02} & Tính giá ngày cuối tuần có phụ phí 15\% (Quy tắc R2) & 1. Chọn thuê trọn T7 và CN. 2. Chọn điểm giao xe cách 4km. 3. Kiểm tra chi tiết giá. & - Tiền thuê: $2 \times 1.150.000 = 2.300.000$đ. - Phí bảo hiểm (10\%): 230.000đ. - Phí giao nhận xe: 0đ (miễn phí $\le 5$km). - Tổng tiền: 2.530.000đ. - Tiền cọc 30\%: 759.000đ. \\ \midrule
\textbf{TC-PRIC-03} & Tính phí giao xe vượt mốc 5km (Quy tắc R3) & 1. Nhập địa chỉ giao xe cách 12km. 2. Quan sát dòng "Phí giao nhận xe". & Phí giao xe tính chuẩn: $(12 - 5) \times 15.000 = 105.000$đ. Tổng tiền = 1.000.000đ (thuê) + 100.000đ (bảo hiểm) + 105.000đ (giao xe) = 1.205.000đ. \\ \midrule
\textbf{TC-PRIC-04} & Kiểm tra giá trị biên khoảng cách giao xe: Đúng 5.0 km & 1. Chọn điểm giao xe đúng 5.0km. & Phí giao nhận xe hiển thị: 0đ (Miễn phí trong 5km). \\ \midrule
\textbf{TC-PRIC-05} & Kiểm tra giá trị biên khoảng cách giao xe: 5.1 km & 1. Chọn điểm giao xe 5.1km. & Phí giao nhận xe tính cho phần vượt: $(5.1 - 5) \times 15.000 = 1.500$đ (hoặc làm tròn theo block 1km tùy cấu hình chủ xe). \\ \midrule
\textbf{TC-PRIC-06} & Áp dụng mã giảm giá phần trăm hợp lệ (Quy tắc R4) & 1. Nhập mã `MIOTOHELLO`. 2. Bấm "Áp dụng". & Hiển thị thông báo "Áp dụng mã thành công". Dòng khuyến mãi ghi nhận: -200.000đ (10\% của 2tr). Tổng tiền giảm trừ chính xác 200.000đ. \\ \midrule
\textbf{TC-PRIC-07} & Áp dụng mã giảm giá cố định (Quy tắc R5) & 1. Nhập mã `BANMOI100`. 2. Bấm "Áp dụng". & Giảm trừ trực tiếp 100.000đ vào tổng thanh toán. Tiền đặt cọc 30\% tính lại theo tổng tiền mới. \\ \midrule
\textbf{TC-PRIC-08} & Áp dụng mã giảm giá đã hết hạn sử dụng (Quy tắc R6) & 1. Nhập mã đã quá hạn. 2. Bấm "Áp dụng". & Báo lỗi: "Mã khuyến mãi này đã hết hạn sử dụng". Không giảm trừ tiền. \\ \midrule
\textbf{TC-PRIC-09} & Áp dụng mã giảm giá khi đơn hàng chưa đạt giá trị tối thiểu & 1. Nhập mã `VIP500`. 2. Bấm "Áp dụng". & Báo lỗi: "Đơn hàng chưa đạt giá trị tối thiểu 1.500.000đ để sử dụng mã này". \\ \midrule
\textbf{TC-PRIC-10} & Áp dụng đồng thời 2 mã giảm giá trên một đơn hàng & 1. Sau khi áp dụng mã 1, cố gắng nhập tiếp mã 2. & Hệ thống thông báo: "Mỗi chuyến đi chỉ được áp dụng tối đa 1 mã khuyến mãi. Mã mới sẽ thay thế mã cũ". \\ \midrule
\textbf{TC-BOOK-01} & Tạo yêu cầu đặt xe thành công (Quy trình thông thường) & 1. Tại trang chi tiết xe, chọn ngày. 2. Nhập ghi chú cho chủ xe: "Tôi đi công tác Vũng Tàu". 3. Bấm "Gửi yêu cầu đặt xe". & - Tạo thành công bản ghi chuyến đi với trạng thái "Chờ chủ xe duyệt". - Hiển thị đồng hồ đếm ngược 60 phút. - Gửi thông báo Push Notification/SMS đến chủ xe. \\ \midrule
\textbf{TC-BOOK-02} & Tạo yêu cầu đặt xe khi chưa được duyệt GPLX & 1. Bấm "Gửi yêu cầu đặt xe". & Hệ thống bật popup cảnh báo: "Bạn cần hoàn tất xác thực Giấy phép lái xe trước khi tiến hành đặt xe", kèm nút chuyển hướng đến trang KYC. \\ \midrule
\textbf{TC-BOOK-03} & Chủ xe phê duyệt yêu cầu trong vòng 60 phút & 1. Chủ xe mở ứng dụng/web, vào danh sách yêu cầu. 2. Bấm nút "Chấp nhận". & Đơn chuyển sang trạng thái "Đã duyệt - Chờ thanh toán cọc". Khách thuê nhận thông báo và link thanh toán cọc trong 60 phút. \\ \midrule
\textbf{TC-BOOK-04} & Chủ xe từ chối yêu cầu đặt xe kèm lý do & 1. Chủ xe bấm "Từ chối". 2. Chọn lý do từ danh sách. 3. Bấm xác nhận. & Đơn chuyển trạng thái "Bị từ chối". Gửi thông báo đến khách thuê kèm gợi ý các mẫu xe tương đương lân cận. Lịch xe được mở lại. \\ \midrule
\textbf{TC-BOOK-05} & Tự động hủy đơn khi chủ xe không phản hồi sau 60 phút & 1. Giả lập thời gian trôi qua quá 60 phút không có thao tác từ chủ xe. & Hệ thống chạy background job chuyển trạng thái đơn sang "Hết hạn phản hồi" (Expired). Gửi thông báo xin lỗi khách hàng. \\ \midrule
\textbf{TC-BOOK-06} & Luồng "Đặt xe nhanh" (Instant Booking) tự động duyệt & 1. Chọn xe có nhãn "Đặt xe nhanh". 2. Bấm "Đặt xe nhanh". & Hệ thống bỏ qua bước chờ chủ xe duyệt, chuyển thẳng đến màn hình thanh toán tiền cọc. \\ \midrule
\textbf{TC-BOOK-07} & Thanh toán cọc thành công qua Thẻ Quốc tế (Visa/MasterCard qua Alepay) & 1. Chọn phương thức "Thẻ tín dụng/ghi nợ Quốc tế". 2. Nhập thông tin thẻ và OTP 3D-Secure. 3. Bấm xác nhận thanh toán. & Cổng thanh toán trả về mã giao dịch thành công. Đơn chuyển trạng thái "Đã đặt cọc" (Booked). Khóa lịch xe trên hệ thống. \\ \midrule
\textbf{TC-BOOK-08} & Thanh toán cọc thất bại do Thẻ không đủ số dư & 1. Nhập thông tin thẻ có số dư 0đ. 2. Xác nhận thanh toán. & Cổng thanh toán báo lỗi: "Giao dịch không thành công do số dư thẻ không đủ". Hệ thống không trừ tiền, giữ nguyên trạng thái đơn để khách chọn thẻ khác. \\ \midrule
\textbf{TC-BOOK-09} & Thanh toán cọc qua Ví điện tử MoMo & 1. Chọn thanh toán "Ví MoMo". 2. Quét mã QR MoMo trên màn hình bằng App MoMo. 3. Xác nhận trừ tiền. & Webhook MoMo IPN gửi tín hiệu thanh toán thành công. Màn hình tự động cập nhật sang trạng thái "Đã đặt cọc thành công". \\ \midrule
\textbf{TC-BOOK-10} & Thanh toán cọc bằng Số dư Ví Mioto (Mioto Wallet) & 1. Chọn phương thức "Ví Mioto". 2. Nhập mã PIN ví. 3. Bấm "Thanh toán". & Khấu trừ 800.000đ trong ví Mioto, số dư mới còn 700.000đ. Đơn chuyển trạng thái "Đã đặt cọc" ngay lập tức mà không qua cổng trung gian. \\ \midrule
\textbf{TC-BOOK-11} & Người dùng hủy ngang khi đang ở trang thanh toán cổng trung gian & 1. Bấm nút hủy trên cổng Alepay. & Hệ thống nhận callback mã hủy, hiển thị thông báo: "Giao dịch chưa hoàn tất. Bạn vẫn còn 45 phút để hoàn thành việc đặt cọc". \\ \midrule
\textbf{TC-TRIP-01} & Thực hiện Check-in bàn giao xe thành công với đủ ảnh & 1. Hai bên gặp mặt, chủ xe/khách mở mục Check-in. 2. Chụp đủ 4 góc ngoại thất và đồng hồ xe. 3. Nhập số ODO ban đầu: 25000. 4. Chọn mức nhiên liệu: Đầy bình. 5. Hai bên ký xác nhận trên màn hình cảm ứng. & Hệ thống lưu trữ toàn bộ ảnh lên Cloud Storage, tạo biên bản bàn giao điện tử. Trạng thái chuyến đi đổi thành "Đang diễn ra" (In Progress). \\ \midrule
\textbf{TC-TRIP-02} & Check-in thiếu ảnh ngoại thất bắt buộc (< 4 ảnh) & 1. Chụp 2 ảnh. 2. Cố gắng bấm nút "Ký xác nhận bàn giao". & Nút bấm bị mờ hoặc báo lỗi: "Vui lòng chụp đầy đủ ít nhất 4 góc xe (đầu, đuôi, sườn trái, sườn phải) để bảo vệ quyền lợi". \\ \midrule
\textbf{TC-TRIP-03} & Check-out trả xe đúng giờ, không vượt km, xăng hoàn trả đủ & 1. Chụp ảnh xe lúc trả. 2. Nhập ODO kết thúc: 25500. 3. Xác nhận xăng đầy bình. 4. Bấm "Hoàn thành chuyến đi". & Hệ thống ghi nhận số km di chuyển là 500km ($\le 600$km). Phụ phí vượt km: 0đ. Phụ phí bù xăng: 0đ. Khách thanh toán 70\% tiền thuê còn lại. \\ \midrule
\textbf{TC-TRIP-04} & Check-out có vượt hạn mức quãng đường (Over-mileage) & 1. Nhập ODO kết thúc: 20450. 2. Hệ thống tự động tính chênh lệch. & Hệ thống tính toán chính xác phần vượt: $(450 - 300) \times 3.000 = 450.000$đ. Tự động thêm dòng "Phụ phí vượt quãng đường: 450.000đ" vào hóa đơn quyết toán. \\ \midrule
\textbf{TC-TRIP-05} & Check-out thiếu xăng so với mức ban đầu & 1. Chọn mức xăng khi trả: 4/8 vạch. 2. Xác nhận tình trạng. & Hệ thống tính phụ phí bù nhiên liệu: $25 \times 24.000 = 600.000$đ cộng vào hóa đơn thanh toán cuối cùng. \\ \midrule
\textbf{TC-TRIP-06} & Check-out trả xe trễ hạn (Over-time surcharge) & 1. Bấm check-out lúc 20:30. 2. Hệ thống so sánh với giờ hẹn trên hợp đồng. & Hệ thống phát hiện trễ 2.5 giờ, tự động tính phụ phí quá giờ: $2.5 \times 100.000 = 250.000$đ và hiển thị trên bảng kê quyết toán. \\ \midrule
\textbf{TC-TRIP-07} & Áp dụng phụ phí vệ sinh xe và phụ phí khử mùi thuốc lá & 1. Chủ xe chọn tích vào "Xe cần vệ sinh" và "Xe có mùi thuốc lá". 2. Tải ảnh minh chứng. 3. Gửi yêu cầu phụ phí. & Hệ thống cộng thêm 400.000đ vào chi phí phát sinh và yêu cầu khách xác nhận thanh toán. \\ \midrule
\textbf{TC-CANC-01} & Khách hủy chuyến trước > 7 ngày (Quy tắc CR-01) & 1. Khách vào chi tiết chuyến đi. 2. Bấm "Hủy chuyến". 3. Chọn lý do hủy. 4. Xác nhận hủy. & - Trạng thái đổi thành "Đã hủy bởi khách". - Thông báo: "Bạn được hoàn lại 100\% tiền cọc (1.000.000đ)". - Tiền được hoàn về Ví Mioto hoặc thẻ gốc trong 3-5 ngày làm việc. - Mở lại lịch trống cho xe. \\ \midrule
\textbf{TC-CANC-02} & Khách hủy chuyến cách giờ nhận xe 3 ngày (Quy tắc CR-02) & 1. Bấm "Hủy chuyến". 2. Xem chính sách bồi cọc hiển thị trên popup. 3. Bấm xác nhận hủy. & - Trạng thái đổi thành "Đã hủy bởi khách". - Khách được hoàn 50\% tiền cọc (500.000đ). - 500.000đ còn lại được chuyển vào ví chủ xe làm tiền bồi thường giữ lịch. \\ \midrule
\textbf{TC-CANC-03} & Khách hủy chuyến gấp dưới 24 giờ trước giờ nhận xe (Quy tắc CR-03) & 1. Bấm "Hủy chuyến". 2. Xác nhận hủy gấp. & - Thông báo: "Theo chính sách, bạn sẽ mất 100\% tiền cọc (1.000.000đ) do hủy chuyến trong vòng 24 giờ". - Không hoàn tiền cọc. - Chuyển tiền đền bù cho chủ xe. \\ \midrule
\textbf{TC-CANC-04} & Kiểm tra giá trị biên thời gian hủy: Đúng 24 giờ 00 phút & 1. Bấm hủy tại mốc đúng 24h00m. & Hệ thống áp dụng chính sách hoàn 50\% cọc (thuộc biên trên của khung 7 ngày - 24h). \\ \midrule
\textbf{TC-CANC-05} & Chủ xe bấm hủy chuyến sau khi khách đã cọc (Quy tắc CR-04) & 1. Chủ xe mở đơn, bấm "Hủy chuyến". 2. Nhập lý do. 3. Xác nhận. & - Khách được hoàn lại 100\% tiền cọc (1.000.000đ). - Khách được tặng thêm voucher đền bù 300.000đ (30\% cọc). - Chủ xe bị phạt 300.000đ khấu trừ vào ví và bị ghi nhận 1 lần hủy chuyến làm giảm tỷ lệ phản hồi. \\ \midrule
\textbf{TC-RATE-01} & Gửi đánh giá 5 sao kèm bình luận và hình ảnh sau khi hoàn tất chuyến & 1. Vào lịch sử chuyến đi, chọn chuyến vừa xong. 2. Bấm "Đánh giá chuyến đi". 3. Chấm 5 sao, nhập bình luận, đính kèm ảnh. 4. Bấm "Gửi đánh giá". & Đánh giá được lưu thành công, hiển thị công khai trên trang chi tiết của xe. Điểm sao trung bình của xe được tính toán lại ngay lập tức. \\ \midrule
\textbf{TC-RATE-02} & Cố gắng gửi đánh giá khi chuyến đi chưa hoàn thành & 1. Mở chi tiết đơn hàng. 2. Tìm nút đánh giá. & Nút "Đánh giá" không xuất hiện hoặc bị vô hiệu hóa kèm chú thích: "Chức năng đánh giá chỉ mở sau khi bạn hoàn tất trả xe". \\ \midrule
\textbf{TC-RATE-03} & Đánh giá lại lần thứ hai cho cùng một chuyến đi & 1. Mở lại chuyến đi đã được đánh giá. 2. Thử bấm gửi đánh giá mới. & Nút đánh giá chuyển thành "Đã đánh giá" và không cho phép gửi thêm lượt đánh giá thứ hai (mỗi chuyến chỉ được đánh giá duy nhất 1 lần). \\ \midrule
\textbf{TC-RATE-04} & Gửi đánh giá sau thời hạn 7 ngày kể từ khi kết thúc chuyến & 1. Truy cập đơn hàng cũ. 2. Kiểm tra quyền đánh giá. & Hệ thống đóng form đánh giá với thông báo: "Thời hạn gửi đánh giá cho chuyến đi này đã kết thúc (trong vòng 7 ngày kể từ khi trả xe)". \\ \midrule
\textbf{TC-RATE-05} & Nhập nội dung bình luận rỗng khi chỉ chấm sao & 1. Chọn 4 sao. 2. Không nhập text. 3. Bấm "Gửi đánh giá". & Hệ thống vẫn cho phép gửi thành công (nội dung văn bản là tùy chọn, điểm sao là bắt buộc). \\ \midrule
\textbf{TC-RATE-06} & Tính toán cập nhật điểm sao trung bình của xe & 1. Gửi đánh giá 5 sao. 2. Tải lại trang chi tiết xe. & Điểm trung bình mới được tính chính xác: $(5 + 5 + 4 + 4 + 5) / 5 = 4.6$ sao. Số lượt đánh giá tăng lên 5 lượt. \\ \midrule
\end{longtable}
